% ── 다리 초안: reynier2003 → 우리 ΔS_vib 음 baseline(§2.3 eq:Svib_mode) ──
% 삽입 위치: ch2_sec03_vibel.tex, vib 문단(현행 33행 jpcc2021 문장) 직후,
%   "전이 폭 안에서 $\Delta S_\vib$..."(현행 35행) 앞.
% 앵커 문장: "...흑연에서는 고-$x$... \emph{음의 baseline}을 만든다(Reynier 의 ``second contribution'')
%   \cite{reynier2003}."(현행 30–31행)
% 컨테이너: srcbox(간결). 우리 기호로 서술.
% ★확인 필요 없음(논문 식 재현 안 함) — 방법 요지는 논문 제목/주제로 검증. 부호 앵커 대응만.

\begin{srcbox}
\textbf{다리 --- $\Delta S_\vib$ 음 baseline 이 어느 실측 앵커인가.} 식~\eqref{eq:Svib_mode} 가 낳는 부분몰
$\Delta S_\vib=\partial S_\vib/\partial x$ 의 고-$x$ 음 baseline 은 Reynier 등\cite{reynier2003} 의 실측 흑연
엔트로피 분해 중 vibrational ``second contribution'' 에 대응한다: 우리 포논 BE 항 $S_\vib$(식~\eqref{eq:Svib_mode})
$\leftrightarrow$ Reynier 의 측정 $\Delta S(x)$ 를 config 발산 항과 vib 음 baseline 으로 가른 둘째 몫.
\emph{방법 요지} --- Reynier 등은 흑연 삽입 엔트로피$\cdot$엔탈피를 온도 의존 전위 측정으로 실측하고, $\Delta
S(x)$ 를 configurational 발산과 vibrational 음 baseline 으로 갈라 해석한다. \emph{가정 차} --- Reynier 는 부호
$\cdot$추세를 두 기여로 정성 분해하나(전이별 정밀값 아님, tier B), 본문은 vib 항을 포논 BE 분포에서
제일원리적으로 세워(식~\eqref{eq:Svib_mode}) 봉우리 폭 안 $T$-무관 상수로 중심 흡수하므로, Reynier 는 고-$x$
음 baseline 의 부호$\cdot$스케일 앵커로만 쓰고 모드별 정량은 제일원리 포논(\cite{jpcc2021}) 소관으로 넘긴다.
\end{srcbox}
